\chapter{Preguntas frecuentes}

\paragraph{No puedo ingresar} Revise que el usuario esté bien escrito (por ejemplo, \texttt{a.romero}, con punto) o toque su foto. Si olvidó la contraseña, use \boton{¿Olvidaste tu contraseña?}. Si el sistema dice que su cuenta está dada de baja, pida a administración que la reactive.

\paragraph{No me llega el código de recuperación} Revise la carpeta de correo no deseado y que el correo sea el registrado en su cuenta. El código vale 10 minutos; si venció, pida otro. Si nunca llegan correos, falta configurar el servicio de correo en el sistema.

\paragraph{No veo una pantalla del menú} Su cargo no la tiene marcada. Pida a quien administra los cargos que la agregue; la verá la próxima vez que inicie sesión.

\paragraph{El sistema me sacó y pide ingresar de nuevo} La sesión dura un turno (8 horas). Ingrese otra vez: vuelve a la pantalla en la que estaba.

\paragraph{El botón Cobrar está gris} Todavía falta algo de cocina. Abra \boton{Marcar listos} y marque lo que ya salió; cuando la columna \textbf{Faltan} quede vacía, se habilita el cobro.

\paragraph{Un producto no se puede elegir en Caja} Está marcado como agotado. Quien tenga acceso a Productos debe encender su interruptor \textbf{Disponible}. Si directamente no aparece, está dado de baja o su categoría lo está.

\paragraph{Una promoción no aparece en Caja} Solo aparecen las vigentes en ese momento. Revise en \pantalla{Promociones} sus fechas, días y horario: la columna Vigencia dice por qué no aplica.

\paragraph{Apareció «Revise el inventario» al enviar un pedido} Algún insumo de la receta no alcanzó: el pedido se envió igual y ese insumo quedó en cero. Cargue la cantidad real en \pantalla{Inventario}.

\paragraph{Me equivoqué en un pedido ya enviado} Lo enviado a cocina no se borra desde Caja. Avise a cocina y, si hace falta, registre lo correcto en un nuevo envío.

\paragraph{El ticket pide usuario y contraseña} Ya se imprimió 3 veces. Un empleado con acceso a Caja debe autorizar la nueva impresión; queda registrado quién lo hizo.

\paragraph{La comanda o el ticket no salen por la impresora} En la aplicación de escritorio, revise en \boton{Impresoras} que estén elegidas la de cocina y la de caja, y pruebe con \boton{Imprimir prueba}. Revise que la impresora esté encendida y con papel. Si la impresión directa falla, use \boton{Imprimir con diálogo}. En el navegador, elija la ticketera en el diálogo de impresión.

\paragraph{El reporte no llega por correo} Revise que el correo esté bien escrito y la carpeta de correo no deseado. Si los correos nunca salen, falta configurar la clave del servicio de correo en esa computadora. En los envíos programados que fallaron, el sistema vuelve a intentar solo a los 15 minutos.

\paragraph{La computadora estaba apagada a la hora de un envío programado} El reporte sale apenas se enciende el sistema, con las ventas desde el envío anterior, y en la lista figura como \textit{Enviado al encender}.

\paragraph{No puedo dar de baja una sección o un cargo} Una sección con mesas ocupadas no se da de baja hasta cobrarlas. Un cargo que tienen empleados activos tampoco: primero cámbieles el cargo.

\paragraph{Borré algo en la eliminación física} No se puede deshacer. Por eso la pantalla muestra antes todo lo que se va y pide escribir \texttt{ELIMINAR}. Si solo quería que algo dejara de usarse, alcanza con darlo de baja.
